%! Trigonometry and Polar Coordinates — Unit 3, Lesson 3.13 — Student Packet Cover Sheet
\documentclass[10pt]{article}
\usepackage{precalculus-article}
\usepackage{precalculus-boxes}
\usepackage{ltablex}
\keepXColumns

\begin{document}

% Full-bleed navy banner
\begin{tikzpicture}[remember picture, overlay]
  \fill[navy] (current page.north west)
    rectangle ([yshift=-0.9in]current page.north east);
\end{tikzpicture}
\vspace{-0.8in}

\begin{center}
  {\color{white}\textbf{\LARGE AP Precalculus}}\\[4pt]
  {\color{skymid}\large\textbf{Unit 3: Trigonometric and Polar Functions}}\\[6pt]
  {\color{white}\textbf{\large Lesson 3.13 \quad Trigonometry and Polar Coordinates}}
\end{center}

\vspace{0.22in}
\namedateperiod
\vspace{0.18in}

\begin{learningtargetbox}
\small
\begin{enumerate}[leftmargin=1.5em, itemsep=5pt]
  \item I can locate a point in the plane using polar coordinates $(r,\theta)$,
    explain what $r$ and $\theta$ represent, and describe how the same point has
    multiple polar representations. \hfill\textit{(LO 3.13.A, EK 3.13.A.1)}
  \item I can convert between polar coordinates $(r,\theta)$ and rectangular
    coordinates $(x,y)$ using $x = r\cos\theta$, $y = r\sin\theta$,
    $r = \sqrt{x^2+y^2}$, and $\theta = \arctan(y/x)$ with a quadrant check.
    \hfill\textit{(LO 3.13.A, EK 3.13.A.2--3)}
  \item I can write a complex number $a+bi$ in polar form $(r\cos\theta)+i(r\sin\theta)$
    by finding the modulus $r$ and argument $\theta$ from the complex plane.
    \hfill\textit{(LO 3.13.A, EK 3.13.A.4)}
\end{enumerate}
\end{learningtargetbox}

\vspace{0.14in}

\begin{tocbox}
\small
\renewcommand{\arraystretch}{1.5}
\begin{tabularx}{\linewidth}{@{} c l X r @{}}
\rowcolor{sky}
\textbf{\#} & \textbf{Component} & \textbf{Description} & \textbf{Score} \\
\midrule
1 & Warm-Up        & Spiral review: distance formula, unit-circle values, inverse tangent & \blank{1.2cm} \\
2 & Guided Notes   & Polar coordinate system; polar$\leftrightarrow$rectangular conversions; complex numbers in polar form & \blank{1.2cm} \\
3 & Group Activity & Tiered practice: polar$\to$rectangular (Tier R), rectangular$\to$polar and complex numbers (Tier A), algebraic proofs (Tier E) & \blank{1.2cm} \\
4 & Exit Ticket    & Convert $(4,\pi/3)$ to rectangular; convert $(-3,3)$ to polar; write $2-2i$ in polar form & \blank{1.2cm} \\
5 & Homework       & Polar and rectangular conversions, multiple representations, complex numbers, AP-style MC & \blank{1.2cm} \\
\midrule
  & \multicolumn{2}{r}{\textbf{Total}} & \blank{1.2cm} \\
\end{tabularx}
\end{tocbox}

\end{document}
